\item La \textit{compresibilidad} $\kappa$ de una sustancia se define como el cambio fraccionario en volumen de dicha sustancia para un cambio dado en presión
$$ \kappa = -\frac{1}{V} \frac{dV}{dP} $$
(a) Explique por qué el signo negativo en esta expresión asegura que $\kappa$ siempre es positiva. (b) Demuestre que si un gas ideal se comprime isotérmicamente, su compresibilidad está dada por $\kappa_1 = 1/P$. (c) \textbf{¿Qué pasaría si?} Demuestre que si un gas ideal se comprime adiabáticamente, su compresibilidad está dada por $\kappa_2 = 1/(\gamma P)$. Determine valores para (d) $\kappa_1$ y (e) $\kappa_2$ para un gas ideal monoatómico a una presión de 2.00 atm.
